\documentclass[9pt]{IEEEtran}

\usepackage[english]{babel}
\usepackage{subcaption}
\usepackage{graphicx}
\usepackage{epstopdf}
\usepackage{fancyhdr}
\usepackage{amsmath}
\usepackage{amsthm}
\usepackage{amssymb}
\usepackage{url}
\usepackage{array}
\usepackage{textcomp}
\usepackage{listings}
\usepackage{hyperref}
\usepackage{xcolor}
\usepackage{colortbl}
\usepackage{float}
\usepackage{gensymb}
\usepackage{longtable}
\usepackage{supertabular}
\usepackage{multicol}

\usepackage[utf8x]{inputenc}

\usepackage[T1]{fontenc}
\usepackage{lmodern}
\input{glyphtounicode}
\pdfgentounicode=1

\graphicspath{{./figures/}}
\DeclareGraphicsExtensions{.pdf,.png,.jpg,.eps}

% correct bad hyphenation here
\hyphenation{op-tical net-works semi-conduc-tor trig-gs}
% ============================================================================================

\title{\vspace{0ex}
Exercise 3: Correlation filter tracking}

\author{Katarina Brešar\vspace{-4.0ex}}

% ============================================================================================

\begin{document}

\maketitle

\section{Introduction}

In this assignment we implemented a correlation filter tracker and evaluated its performance on the VOT2014 dataset.
The tracker operates by learning a filter in the Fourier domain that produces a Gaussian peak response at the
target's location. The filter is then updated over time using a learning rate to adapt to changes in the target's appearance and 
location. We then evaluated the tracker's performance using VOT-lite toolkit that uses standard performance measures
accuracy and robustness.

\section{Experiments}
The tracker was implemented using a simplified MOSS approach. 

\subsection{Analysis of parameters}
We analyzed how the learning rate $\alpha$ and the gaussian kernel bandwidth $\sigma$ influence the tracker's performance.
Testing the tracker on the different sequences from the VOT2014 dataset, we found that a high learning rate (close to $1$) allows the tracker
to quickly adapt to changes in the targets appearance, but it can lead to instability, as it makes it more sensitive to occlusions (TODO woman example).
Our thought was that a higher learning rate would then work best on the sequence with more dynamic targets (like in bolt and gymnastic). However, when
tested on such sequences, we found that for all a learning rate at around $0.1$ to $0.3$ gives the best performance, when it comes to
average overlap between the predicted and the ground truth bounding boxes. (TODO: add plots for all three sequences of the avg overlap).
As seen in figure (TODO add) the graph is reverse U shaped, with the best performance close to the middle.

Focusing on the kernel bandwidth $\sigma$


\subsection{Larger patch}
With the idea of capturing more background for better filter learning, we tested the tracker with patches larger than the target.
So we extracted patches that were enlarged by some factor and then trained the filter on those. But the reported region (so the bounding box that the tracker returns) was
still the same as the target size. Just to calculate the new center we used a larger window. This approach sometimes improved the performance,
TODO especially when the target was moving fast and changing its appearance (gymnastics?). But a larger patch
also made the tracker more sensitive to occlusions and background clutter, which could lead to more tracking failures. 
TODO add example of failures with larger patches. 






\subsection{Tracking speed}
Using the VOT-lite toolkit, we tested our tracker on the VOT2014 dataset and got the results seen in table~\ref{tab:tracker_results}.


\begin{table}[h!]
    \centering
    \resizebox{0.25\textwidth}{!}{
    \begin{tabular}{|c|c|}
    \hline
    Average overlap  & 0.47 \\ 
    Total failures & 66 \\
    Average speed (FPS) & 814.30 \\
    Total frames & 10213 \\ \hline
    \end{tabular}}
    \caption{Tracker results on the VOT2014 dataset.}\label{tab:tracker_results}
\end{table}


\section{Conclusion}



\bibliographystyle{IEEEtran}
\bibliography{bibliography}

\end{document}
